% ============================================================================
%  scs.tex -- SCS: the method and the implementation
%
%  A standalone write-up of the algorithm SCS implements and of the software
%  in this repository. It is not generated from the Sphinx sources in
%  docs/src; it is a hand-maintained companion to them, and the two should be
%  kept consistent when the algorithm changes.
%
%  Build:  make          (in this directory; needs pdflatex)
%          make clean
% ============================================================================
\documentclass[11pt,a4paper]{article}

\usepackage[margin=1in]{geometry}
\usepackage{amsmath,amssymb,amsthm}
\usepackage{booktabs}
\usepackage{array}
\usepackage[colorlinks=true,linkcolor=blue,citecolor=blue,urlcolor=blue]{hyperref}

% --- notation ---------------------------------------------------------------
\newcommand{\R}{\mathbf{R}}
\newcommand{\Sym}{\mathbf{S}}
\newcommand{\K}{\mathcal{K}}
\newcommand{\C}{\mathcal{C}}
\newcommand{\Q}{\mathcal{Q}}
\newcommand{\N}{\mathcal{N}}
\newcommand{\proj}[1]{\Pi_{#1}}
\newcommand{\tr}{^\top}
\DeclareMathOperator*{\minimize}{minimize}
\DeclareMathOperator*{\maximize}{maximize}
\DeclareMathOperator{\diag}{\mathbf{diag}}
\DeclareMathOperator{\dist}{\mathbf{dist}}

\theoremstyle{plain}
\newtheorem{proposition}{Proposition}

% Authorship: SCS is a community project whose primary maintainer is Brendan
% O'Donoghue; the algorithm is due to the authors cited in Section 10. Replace
% the line below with the author list appropriate to wherever this document is
% used.
\title{\textbf{SCS: a splitting conic solver}\\[0.4em]
       \large The method, and the software that implements it}
\author{The SCS developers}
\date{\today}

\begin{document}
\maketitle

\begin{abstract}
SCS solves large-scale convex quadratic cone programs. It embeds the primal-dual
optimality conditions of such a program into a homogeneous self-dual feasibility
problem, and applies Douglas--Rachford splitting to that embedding under a
diagonal metric. Each iteration costs one solve of a fixed quasi-definite linear
system and one Euclidean projection onto a product of cones; because the linear
system does not change between iterations it can be factorized once and reused,
or solved approximately by an iterative method. The iterates satisfy the conic
constraints and complementary slackness exactly at every step, so termination
reduces to checking three residuals, and the same iterates yield a certificate
of primal or dual infeasibility when no solution exists. This document derives
the method, describes the heuristics that make it work in practice --- data
equilibration, an adaptive diagonal metric, and Anderson acceleration --- and
then documents the software in this repository: its layout, its C API, its
pluggable linear-solver backends, its build configuration and its tests.
\end{abstract}

\tableofcontents

% ===========================================================================
\section{Introduction}
% ===========================================================================

SCS (\emph{splitting conic solver}) is a numerical optimization package for
convex cone programs with a quadratic objective. It is a first-order method: it
does not form or factorize a Newton system involving the barrier of the cone,
and its per-iteration cost is one linear solve with a \emph{fixed} matrix plus
one cone projection. This buys three properties that distinguish it from
interior-point solvers.

\begin{itemize}
\item \textbf{Scale.} The linear system is the same at every iteration, so a
  single factorization is amortized over the whole solve; alternatively it can
  be solved matrix-free by conjugate gradients, which needs only products with
  $P$, $A$ and $A\tr$.
\item \textbf{Infeasibility detection.} The method works on a homogeneous
  embedding, in which an optimal point and a certificate of infeasibility are
  both fixed points of the same operator. Nothing special has to be done to
  detect an infeasible or unbounded problem.
\item \textbf{Warm starts.} The iterate is a single vector in the embedding
  space, so a solution of a nearby problem is a usable starting point, and the
  factorization and equilibration can be cached across solves.
\end{itemize}

The cost is accuracy: first-order methods converge slowly to high precision.
Sections~\ref{sec:scaling} and~\ref{sec:aa} describe the three mechanisms SCS
uses to fight this --- equilibration of the data, an adaptively updated
diagonal metric, and Anderson acceleration of the fixed-point iteration.

The algorithm in this repository is the one derived in
\cite{odonoghue:21}, which extends the original SCS
method \cite{ocpb:16} to quadratic objectives and replaces its
self-dual embedding with a homogeneous embedding of the linear complementarity
problem. Where the two differ, this document follows the former.

% ===========================================================================
\section{The quadratic cone program}
% ===========================================================================

SCS solves the primal-dual pair
\begin{equation}\label{eq:problem}
\begin{array}{ll}
\minimize & (1/2)x\tr P x + c\tr x\\
\mbox{subject to} & Ax + s = b\\
                  & s \in \K
\end{array}
\qquad\qquad
\begin{array}{ll}
\maximize & -(1/2)x\tr P x - b\tr y\\
\mbox{subject to} & Px + A\tr y + c = 0\\
                  & y \in \K^*
\end{array}
\end{equation}
over the primal variable $x \in \R^n$, the dual variable $y \in \R^m$ and the
slack $s \in \R^m$. The data are a sparse symmetric positive semidefinite
$P \in \Sym_+^n$, a sparse $A \in \R^{m \times n}$, dense vectors
$c \in \R^n$ and $b \in \R^m$, and a nonempty closed convex cone
$\K \subseteq \R^m$ with dual cone $\K^*$.

\subsection{Optimality conditions}

When strong duality holds the Karush--Kuhn--Tucker conditions are necessary and
sufficient for optimality \cite{boyd:04}. For \eqref{eq:problem} they read
\begin{equation}\label{eq:kkt}
Ax + s = b,\qquad Px + A\tr y + c = 0,\qquad
s \in \K,\qquad y \in \K^*,\qquad s \perp y,
\end{equation}
that is, primal feasibility, dual feasibility, primal and dual cone membership,
and complementary slackness. At any point satisfying the first four conditions,
the last is equivalent to a zero duality gap,
\begin{equation}\label{eq:gap}
s \perp y \quad\Longleftrightarrow\quad c\tr x + b\tr y + x\tr P x = 0 .
\end{equation}

\subsection{Certificates of infeasibility}

If no optimal solution exists, \eqref{eq:kkt} has no solution, but one of two
certificates does. Any $y \in \R^m$ with
\begin{equation}\label{eq:pinfeas}
A\tr y = 0,\qquad y \in \K^*,\qquad b\tr y < 0
\end{equation}
proves the primal problem infeasible (equivalently, the dual unbounded); and
any $x \in \R^n$ with
\begin{equation}\label{eq:dinfeas}
Px = 0,\qquad -Ax \in \K,\qquad c\tr x < 0
\end{equation}
proves the dual infeasible (the primal unbounded). Section~\ref{sec:term} gives
the numerical tests SCS actually applies.

\subsection{Cones}\label{sec:cones}

$\K$ is a Cartesian product of the primitive cones in
Table~\ref{tab:cones}. The rows of $A$ and $b$ must be ordered to match the
order of that table, and within each cone must follow the order in the
mathematical description --- for the box and second-order cones, for instance,
the scalar $t$ comes first.

\begin{table}[t]
\centering
\small
\begin{tabular}{@{}l>{$}l<{$}l@{}}
\toprule
Cone & \text{Definition} & \texttt{ScsCone} field \\
\midrule
Zero          & \{s \mid s = 0\}                                          & \texttt{z} \\
Nonnegative   & \{s \mid s \geq 0\}                                       & \texttt{l} \\
Box           & \{(t,s) \in \R \times \R^k \mid t\,l \leq s \leq t\,u\} & \texttt{bl}, \texttt{bu}, \texttt{bsize} \\
Second-order  & \{(t,s) \in \R \times \R^k \mid \|s\|_2 \leq t\}          & \texttt{q}, \texttt{qsize} \\
Semidefinite  & \{s \in \R^{k(k+1)/2} \mid \mathrm{mat}(s) \succeq 0\}    & \texttt{s}, \texttt{ssize} \\
Complex SD    & \{s \in \R^{k^2} \mid \mathrm{mat}(s) \succeq 0\}         & \texttt{cs}, \texttt{cssize} \\
Exponential   & \{(x,y,z) \mid y e^{x/y} \leq z,\ y > 0\}                 & \texttt{ep} \\
Dual exp.     & \{(u,v,w) \mid -u e^{v/u} \leq e w,\ u < 0\}              & \texttt{ed} \\
Power         & \{(x,y,z) \mid x^p y^{1-p} \geq |z|\}                     & \texttt{p}, \texttt{psize} \\
\midrule
\multicolumn{3}{@{}l@{}}{\emph{Spectral cones (experimental; require \texttt{USE\_SPECTRAL\_CONES=1})}}\\
Log-determinant & \{(t,v,X) \mid t \leq v \log\det(X/v),\ v>0,\ X \succ 0\} & \texttt{d}, \texttt{dsize} \\
Nuclear norm    & \{(t,X) \mid \|X\|_* \leq t\}                             & \texttt{nuc\_m}, \texttt{nuc\_n} \\
$\ell_1$ norm   & \{(t,x) \mid \|x\|_1 \leq t\}                             & \texttt{ell1}, \texttt{ell1\_size} \\
Sum-of-largest-$\lambda$ & \{(t,X) \mid \textstyle\sum_{i=1}^k \lambda_i(X) \leq t\} & \texttt{sl\_n}, \texttt{sl\_k} \\
\bottomrule
\end{tabular}
\caption{The primitive cones. Semidefinite cones are passed in a scaled
lower-triangular vectorization; see the API documentation for the exact
convention.}
\label{tab:cones}
\end{table}

All the algorithm needs from a cone is its Euclidean projection: given $z_0$,
solve
\begin{equation}\label{eq:coneproj}
\proj{\K}(z_0) = \operatorname*{argmin}_{z \in \K} \|z - z_0\|_2^2 .
\end{equation}
For every cone in Table~\ref{tab:cones} this has either a closed form or a
cheap one-dimensional root-find; the semidefinite and spectral cones need an
eigendecomposition, which is why they require LAPACK.

% ===========================================================================
\section{The homogeneous embedding}\label{sec:embedding}
% ===========================================================================

The optimality conditions \eqref{eq:kkt} are a monotone complementarity problem
in $(x,y,s)$. To make them homogeneous --- so that infeasibility appears as a
solution rather than as the absence of one --- introduce a scalar
$\tau \geq 0$ that multiplies $b$ and $c$, and a scalar $\kappa \geq 0$
complementary to it.

Let $u = (x, y, \tau) \in \R^{n+m+1}$ and let
\begin{equation}\label{eq:Q}
\Q = \begin{bmatrix} P & A\tr & c \\ -A & 0 & b \\ -c\tr & -b\tr & 0
     \end{bmatrix},
\qquad
\C = \R^n \times \K^* \times \R_+ ,
\end{equation}
so that $\Q$ is the sum of a positive semidefinite and a skew-symmetric matrix,
hence monotone, and $\C$ is a closed convex cone. The embedded problem is
\begin{equation}\label{eq:embedding}
\text{find } u \in \C \quad\text{such that}\quad 0 \in \Q u + \N_\C(u),
\end{equation}
where $\N_\C$ is the normal cone of $\C$. A solution always exists.

Write $v = \Q u = (r, s, \kappa)$. Unpacking \eqref{eq:embedding} block by
block gives $r = Px + A\tr y + c\tau = 0$, then $-Ax + b\tau = s$ with
$s \in \K$, $s \perp y$, and finally $-c\tr x - b\tr y = \kappa \geq 0$ with
$\tau \kappa = 0$. Two cases follow.

\begin{itemize}
\item \textbf{$\tau > 0$, $\kappa = 0$.} Then $(x, y, s)/\tau$ satisfies
  \eqref{eq:kkt} exactly: dividing through, $Ax + s = b$, $Px + A\tr y + c = 0$,
  $s \in \K$, $y \in \K^*$ and $c\tr x + b\tr y = 0$, which by \eqref{eq:gap}
  is complementary slackness. We recover the solution of \eqref{eq:problem} by
  scaling.
\item \textbf{$\tau = 0$, $\kappa > 0$.} Then $c\tr x + b\tr y = -\kappa < 0$
  while $A\tr y = -Px$, $-Ax = s \in \K$, $y \in \K^*$. At least one of
  $b\tr y < 0$ and $c\tr x < 0$ holds, giving a certificate
  \eqref{eq:pinfeas} or \eqref{eq:dinfeas}. If a problem is simultaneously
  primal and dual infeasible, SCS returns whichever certificate it finds first;
  the usual reading of primal infeasibility as dual unboundedness does not
  apply in that case.
\end{itemize}

The pair $(\tau = 0, \kappa = 0)$ is possible but uninformative; it corresponds
to a problem that is neither solvable nor certifiably infeasible in the above
sense, and in practice shows up as the residuals stalling.

% ===========================================================================
\section{Douglas--Rachford splitting}\label{sec:dr}
% ===========================================================================

Problem \eqref{eq:embedding} asks for a zero of the sum of two monotone
operators, $\Q$ (linear, monotone) and $\N_\C$ (the normal cone, maximal
monotone). Douglas--Rachford splitting \cite{lions:79} handles exactly this
sum: from an initial $w^0$, for $k = 0, 1, \dots$,
\begin{align}\label{eq:dr}
\tilde u^{k+1} &= (I + \Q)^{-1} w^k, \nonumber\\
u^{k+1} &= \proj{\C}\!\left(2\tilde u^{k+1} - w^k\right), \\
w^{k+1} &= w^k + u^{k+1} - \tilde u^{k+1}, \nonumber
\end{align}
using that the resolvent of $\N_\C$ is the projection onto $\C$. The iterates
satisfy $u^k \to u^\star$ with $u^\star$ solving \eqref{eq:embedding}, and
$w^k \to u^\star + \Q u^\star$.

\subsection{Relaxation}

The third step admits a relaxation parameter $\alpha \in (0,2)$:
\begin{equation}\label{eq:relax}
w^{k+1} = w^k + \alpha\left(u^{k+1} - \tilde u^{k+1}\right).
\end{equation}
Plain Douglas--Rachford is $\alpha = 1$; $\alpha < 1$ is under-relaxation and
$\alpha > 1$ over-relaxation. Values around $1.5$ work well in practice and
$\alpha = 1.5$ is the default (setting \texttt{alpha}). At $\alpha = 2$ the
method becomes Peaceman--Rachford splitting, which is not convergent in
general.

\subsection{A diagonal metric}\label{sec:metric}

Convergence is much improved by running the splitting in a metric other than
the Euclidean one. Replace $I$ by a positive diagonal
$R \in \R^{(n+m+1)\times(n+m+1)}$:
\begin{align}\label{eq:drR}
\tilde u^{k+1} &= (R + \Q)^{-1} R\, w^k, \nonumber\\
u^{k+1} &= (R + \N_\C)^{-1} R \left(2\tilde u^{k+1} - w^k\right), \\
w^{k+1} &= w^k + \alpha\left(u^{k+1} - \tilde u^{k+1}\right), \nonumber
\end{align}
which converges to $w^k \to u^\star + R^{-1}\Q u^\star$ for the same
$u^\star$. The relaxation $\alpha$ and the metric $R$ do not interact, so the
two modifications compose directly.

The second step of \eqref{eq:drR} is a projection in the $R$-norm
$\|x\|_R = \sqrt{x\tr R x}$, which for a general $R$ is not the Euclidean
projection \eqref{eq:coneproj} and would need a different routine for every
cone. SCS avoids this by constraining $R$ to be constant within each cone
block.

\begin{proposition}\label{prop:metric}
Let $\K$ be a cone and $r > 0$ a scalar. Then
$(rI + \partial I_\K)^{-1}(r y) = \proj{\K}(y)$, where $I_\K$ is the indicator
of $\K$. Consequently, if $R$ is constant on each cone block of $\C$, then
$(R + \partial I_\C)^{-1} R = (I + \partial I_\C)^{-1} = \proj{\C}$.
\end{proposition}

\begin{proof}
Let $x = (rI + \partial I_\K)^{-1}(ry)$. Then
$0 \in r(x - y) + \partial I_\K(x)$, and dividing by $r > 0$ gives
$0 \in (x - y) + \partial I_\K(x)$, whose solution is
$x = \operatorname*{argmin}_{z\in\K}\|z-y\|_2^2$.
\end{proof}

So with $R$ blockwise constant the existing Euclidean projections are reused
unchanged, and $R$ only affects the linear solve. SCS takes
\begin{equation}\label{eq:R}
R = \begin{bmatrix} \rho_x I_n & 0 & 0 \\
                    0 & \diag(\rho_y) & 0 \\
                    0 & 0 & d \end{bmatrix},
\qquad \rho_x \in \R,\ \rho_y \in \R^m,\ d \in \R .
\end{equation}
Here $\rho_x$ is the setting \texttt{rho\_x} (default $10^{-6}$), $d$ is the
constant \texttt{TAU\_FACTOR} in \texttt{include/glbopts.h} (default $10$), and
$\rho_y \approx (1/\texttt{scale})\,I$ with the \texttt{scale} parameter
updated adaptively as described in Section~\ref{sec:adaptive}. The zero cone is
the one exception to the approximation, taking
$\rho_y = 1/(1000\,\texttt{scale})$; the box cone is a mild exception to
blockwise constancy. Note $d$ enters the algorithm only through the scalar
root-find of Section~\ref{sec:linsys}, so there is considerable freedom in
choosing it, and how best to do so is open.

\subsection{Recovering the dual variable}

The iterate $u^k = (x^k, y^k, \tau^k)$ carries the primal and dual variables,
but not the slack. The companion vector
\begin{equation}\label{eq:v}
v^{k+1} = R\left(u^{k+1} + w^k - 2\tilde u^{k+1}\right) \to \Q u^\star
\end{equation}
supplies $(r^k, s^k, \kappa^k)$. That $v^{k+1} \in \C^*$ follows from the
Moreau decomposition under a diagonal metric,
\begin{equation}\label{eq:moreau}
x + R^{-1}\proj{\C^*}^{R^{-1}}(-Rx) = \proj{\C}^{R}(x),
\end{equation}
applied to \eqref{eq:v}:
\[
v^{k+1}
= R\left(\proj{\C}^R(2\tilde u^{k+1} - w^k) + w^k - 2\tilde u^{k+1}\right)
= \proj{\C^*}^{R^{-1}}\!\left(R(w^k - 2\tilde u^{k+1})\right) \in \C^* .
\]
The two components of \eqref{eq:moreau} are $R$-orthogonal, so
$v^k \perp u^k$. This is the reason the iterates satisfy $s \in \K$,
$y \in \K^*$ and $s \perp y$ \emph{exactly} at every iteration --- a fact the
termination criteria of Section~\ref{sec:term} rely on.

% ===========================================================================
\section{The two expensive steps}
% ===========================================================================

\subsection{The linear system}\label{sec:linsys}

The first step of \eqref{eq:drR} requires applying $(R + \Q)^{-1}$, a system of
size $n + m + 1$ whose last row and column carry the dense vectors $b$ and $c$.
SCS reduces it to a system of size $n + m$ with a fixed sparse matrix, plus a
scalar quadratic.

Split $w^k = (\mu^k, \eta^k)$ with $\mu^k \in \R^{n+m}$ and
$\eta^k \in \R$, write $R_{-1}$ for the leading $(n+m)$ block of $R$, and set
\begin{equation}
M = \begin{bmatrix} P & A\tr \\ -A & 0 \end{bmatrix},
\qquad
h = \begin{bmatrix} c \\ -b \end{bmatrix} .
\end{equation}
Then $\tilde u^{k+1}_{-1} = p^k - \tau^{k+1} g$, where
\begin{equation}\label{eq:pg}
p^k = (R_{-1} + M)^{-1} R_{-1}\mu^k,
\qquad
g = (R_{-1} + M)^{-1} h ,
\end{equation}
and $\tau^{k+1}$ is the positive root of $a\tau^2 + b\tau + c_0 = 0$ with
\begin{equation}\label{eq:rootplus}
a = d + g\tr R_{-1} g, \qquad
b = \mu^{k\top} R_{-1} g - 2 p^{k\top} R_{-1} g - d\,\eta^k, \qquad
c_0 = p^{k\top} R_{-1} p^k - p^{k\top} R_{-1}\mu^k .
\end{equation}
The vector $g$ depends only on the data and on $R$, so it is computed once
(and recomputed whenever $R$ changes) and cached; each iteration therefore
costs one solve for $p^k$, five weighted inner products, and a root-find. This
is the \texttt{root\_plus} routine in \texttt{src/scs.c}.

The remaining solve is with
\begin{equation}\label{eq:kkt-mat}
\begin{bmatrix} R_x + P & A\tr \\ A & -R_y \end{bmatrix}
\begin{bmatrix} x \\ y \end{bmatrix}
=
\begin{bmatrix} z^k_x \\ z^k_y \end{bmatrix},
\end{equation}
the sign of the $R_y$ block flipped to make the matrix quasi-definite (the
solution of the original system is recovered from this one). Two properties
matter:

\begin{itemize}
\item The matrix is \emph{constant} across iterations, changing only when
  \texttt{scale} is updated. A direct method factorizes it once and reuses the
  factors.
\item The system may be solved \emph{approximately}, provided the errors are
  summable. This is what admits iterative and matrix-free solvers.
\end{itemize}

Table~\ref{tab:linsys} lists the backends in \texttt{linsys/}. The default
sparse direct solver computes a permuted $LDL\tr$ factorization using the
bundled AMD \cite{amd} ordering and QDLDL \cite{qdldl}; quasi-definiteness
gives strong existence guarantees for the factorization. The indirect solver
instead eliminates $y$ to obtain the positive definite reduced system
\begin{equation}\label{eq:reduced}
\left(R_x + P + A\tr R_y^{-1} A\right) x = z^k_x + A\tr R_y^{-1} z^k_y,
\qquad y = R_y^{-1}\left(Ax - z^k_y\right),
\end{equation}
and solves it by conjugate gradients, each iteration costing one product each
with $P$, $A$ and $A\tr$. The tolerance of each solve is a small fraction
(\texttt{CG\_TOL\_FACTOR}) of the smaller of the current primal and dual
residuals, capped by a term decaying like $O(1/k^{\texttt{CG\_RATE}})$ and
floored at \texttt{CG\_BEST\_TOL}; solves are warm-started from the previous
solution. The implementation also harvests approximate eigenvectors of the
reduced matrix from the CG iterations themselves (eigCG, \cite{eigcg}) and
deflates the smallest eigenvalues from later solves, which cuts CG iterations
substantially on ill-conditioned problems at no extra matrix-multiply cost.

\begin{table}[t]
\centering
\small
\begin{tabular}{@{}llp{7.2cm}@{}}
\toprule
Backend & Directory & Method \\
\midrule
Sparse direct (default) & \texttt{linsys/cpu/direct}     & Permuted $LDL\tr$ via AMD + QDLDL, factors cached \\
Sparse indirect         & \texttt{linsys/cpu/indirect}   & Conjugate gradients on \eqref{eq:reduced}, with eigCG deflation \\
Dense direct            & \texttt{linsys/cpu/dense}      & Cholesky (\texttt{dpotrf}) of the $n \times n$ Gram matrix $R_x + P + A\tr R_y^{-1}A$ \\
MKL Pardiso             & \texttt{linsys/mkl/direct}     & Intel oneMKL parallel sparse direct solver \\
Apple Accelerate        & \texttt{linsys/accelerate/direct} & Accelerate unpivoted sparse $LDL\tr$ (macOS only; no \texttt{DLONG}) \\
cuDSS                   & \texttt{linsys/cudss/direct}   & NVIDIA cuDSS direct solve on the GPU \\
GPU indirect            & \texttt{linsys/gpu/indirect}   & Conjugate gradients with GPU matrix products \\
\bottomrule
\end{tabular}
\caption{Linear-solver backends. Each is built into its own binary; a new one
is added by implementing \texttt{ScsLinSysWork} and the functions declared in
\texttt{include/linsys.h}.}
\label{tab:linsys}
\end{table}

\subsection{The cone projection}

The second step of \eqref{eq:drR} is the projection \eqref{eq:coneproj} onto
$\C = \R^n \times \K^* \times \R_+$, which splits into independent projections
onto the blocks of Table~\ref{tab:cones}. Most are closed-form. The
semidefinite blocks require an eigendecomposition of $\mathrm{mat}(s)$ and
clipping of the negative eigenvalues; the exponential cone needs a damped
Newton root-find, with a bisection fallback if it fails to converge
(\texttt{src/exp\_cone.c}); the spectral cones
(\texttt{src/spectral\_cones}) reduce to problems on the eigenvalues or
singular values of the matrix argument. Because $\C$ is a product, this step
parallelizes across blocks, which is what \texttt{USE\_OPENMP} exploits.

% ===========================================================================
\section{Scaling}\label{sec:scaling}
% ===========================================================================

The practical performance of a first-order method depends heavily on how the
problem is scaled. SCS attacks this in two places: once, statically, on the
data, and repeatedly, dynamically, on the metric $R$.

\subsection{Data equilibration}

With the \texttt{normalize} setting enabled (the default), SCS computes
diagonal matrices $E \in \R^{n \times n}$, $D \in \R^{m \times m}$ and a scalar
$\sigma > 0$ that whiten the data,
\begin{equation}\label{eq:equil}
\hat P = EPE, \qquad \hat A = DAE, \qquad \hat c = \sigma E c, \qquad
\hat b = \sigma D b,
\end{equation}
which is precisely a symmetric rescaling of the matrix appearing in the
embedding \eqref{eq:Q}:
\begin{equation}
\begin{bmatrix} \hat P & \hat A\tr & \hat c \\
                \hat A & 0 & \hat b \\
                \hat c\tr & \hat b\tr & 0 \end{bmatrix}
=
\begin{bmatrix} E & & \\ & D & \\ & & \sigma \end{bmatrix}
\begin{bmatrix} P & A\tr & c \\ A & 0 & b \\ c\tr & b\tr & 0 \end{bmatrix}
\begin{bmatrix} E & & \\ & D & \\ & & \sigma \end{bmatrix} .
\end{equation}
To this matrix SCS applies $25$ steps of Ruiz equilibration
\cite{ruiz:01} followed by one step of $\ell_2$ equilibration. Both counts are
compile-time constants (\texttt{NUM\_RUIZ\_PASSES} and
\texttt{NUM\_L2\_PASSES}) in \texttt{linsys/scs\_matrix.c}. The $\ell_2$ pass
uses
root-mean-square norms --- $\ell_2$ norms divided by the number of nonzeros ---
which keeps rows and columns of very different lengths comparable in the
stacked matrix.

The one complication is that $D$ must preserve cone membership,
$s \in \K \Leftrightarrow D_\K s \in \K$ for each sub-cone. As with the metric
$R$, this is enforced by restricting $D_\K = d_\K I_\K$ to be constant within
each cone; the scalar $d_\K$ is the $\ell_\infty$ norm of the computed in-cone
values for the Ruiz passes, and the average of the in-cone root-mean-square
norms for the $\ell_2$ pass.

Equilibration moves the fixed point, so solutions and residuals are mapped
back. From the equilibrated solution $(\hat x, \hat y, \hat s)$,
\begin{equation}\label{eq:unequil}
x = E\hat x/\sigma, \qquad y = D\hat y/\sigma, \qquad s = D^{-1}\hat s/\sigma,
\end{equation}
and the residuals of the original problem are computed from equilibrated
quantities as
\begin{equation}\label{eq:unequil-res}
r_p = \tfrac{1}{\sigma}\left\|D^{-1}(\hat A\hat x + \hat s - \hat b)\right\|,
\quad
r_d = \tfrac{1}{\sigma}\left\|E^{-1}(\hat P\hat x + \hat A\tr\hat y + \hat c)\right\|,
\quad
r_g = \tfrac{1}{\sigma^2}\left|\hat x\tr\hat P\hat x + \hat b\tr\hat y + \hat c\tr\hat x\right| ,
\end{equation}
under any norm. Analogous relations hold for the infeasibility certificate
quantities. Older versions of SCS carried separate \texttt{primal\_scale} and
\texttt{dual\_scale} parameters for $c$ and $b$; both are now the single
$\sigma$.

\subsection{Adaptive metric updates}\label{sec:adaptive}

The $\rho_y$ block of $R$ in \eqref{eq:R} is governed by the \texttt{scale}
parameter, whose best value is strongly problem dependent. With
\texttt{adaptive\_scale} enabled (the default) SCS adjusts it during the solve,
aiming to balance the convergence of the primal and dual residuals: loosely,
\texttt{scale} increases when the primal residual dominates and decreases when
the dual does.

Concretely, let $l$ be the number of iterations since the last update, write
$(x, y, \tau) = u^k$ and $(0, s, \kappa) = v^k$, and define the relative
residuals
\begin{equation}\label{eq:relres}
\hat r^k_p = \frac{\|Ax + s - b\tau\|}{\max(\|Ax\|, \|s\|, \|b\tau\|)},
\qquad
\hat r^k_d = \frac{\|Px + A\tr y + c\tau\|}{\max(\|Px\|, \|A\tr y\|, \|c\tau\|)} ,
\end{equation}
in the $\ell_\infty$ norm by default (\texttt{SCALE\_NORM} in
\texttt{glbopts.h}). Let
\begin{equation}\label{eq:beta}
\beta = \left(\prod_{i=0}^{l-1}
        \frac{\hat r^{k-i}_p}{\hat r^{k-i}_d}\right)^{1/l}
\end{equation}
be the geometric mean of their ratio over those iterations. If $\beta$ leaves a
band around $1$ (say $[1/3, 3]$) and at least \texttt{RESCALING\_MIN\_ITERS}
(default $100$) iterations have passed since the last update, then
\begin{equation}\label{eq:scaleupdate}
\texttt{scale}^+ = \sqrt{\beta}\;\texttt{scale} .
\end{equation}
The square root damps the step, to avoid overshooting and oscillating around
the best value. Updates are not free: with a direct linear solver a new
\texttt{scale} means refactorizing \eqref{eq:kkt-mat}, and because the operator
itself has changed, the Anderson acceleration history must be reset. Hence the
minimum-iteration guard.

\subsection{Per-row refinement}

The update \eqref{eq:scaleupdate} moves the whole $\rho_y$ block together. With
\texttt{adaptive\_diag\_scale} (also on by default) each row additionally
carries its own multiplier, updated at scale-update events from that row's
relative primal residual
\begin{equation}\label{eq:rowres}
\hat r^k_{p,i} = \frac{|(Ax + s - b\tau)_i|}
                      {\max(|(Ax)_i|,\ |s_i|,\ |b_i\tau|,\ \epsilon_{\rm den})} ,
\end{equation}
by a damped multiplicative step towards the geometric mean of the profile,
clamped to a bounded range. This is the dynamic analogue of the row half of
Ruiz equilibration, driven by the residuals the algorithm actually produces
rather than by norms of the data.

Within each non-polyhedral cone block the multipliers are constrained to a form
the projection can absorb --- by Proposition~\ref{prop:metric} a single value
per block suffices in general, but real semidefinite blocks admit the richer
rank-one family $w_{ij} = \delta_i\delta_j$, corresponding to the diagonal
congruence $X \mapsto DXD$ which maps the cone to itself. That family is fitted
per block by least squares in log space, and the projection under the resulting
metric is computed exactly by a congruence sandwich around the
eigendecomposition.

The floor $\epsilon_{\rm den}$ in \eqref{eq:rowres} is a small fraction of the
root-mean-square denominator across rows, and it matters more than it looks. A
row with $b_i = 0$ divides only by $\max(|(Ax)_i|, |s_i|)$, both of which
shrink as the iterates converge, so its ratio inflates even while the row is
converging perfectly well. Without the floor the profile would partly report
which terms a row happens to contain rather than how it is converging, and the
metric would chase that artifact. Rows with denominators of ordinary size are
unaffected. A scale update also fires when the accumulated per-row drift alone
is large enough, so that a \texttt{scale} pinned at its limits cannot freeze
the diagonal.

% ===========================================================================
\section{Anderson acceleration}\label{sec:aa}
% ===========================================================================

The iteration \eqref{eq:drR} is a fixed-point iteration $x^{k+1} = f(x^k)$ for
the map $f$ that takes $w^k$ to $w^{k+1}$. Anderson acceleration (AA) is a
quasi-Newton method for such iterations \cite{anderson:65}: instead of taking
$f(x^k)$, it takes a linear combination of the last $m_k + 1$ outputs of the
map,
\begin{equation}\label{eq:aa}
\begin{array}{l}
\text{for } k = 0, 1, \dots \\
\quad m_k = \min\{m, k\} \\
\quad \text{choose } \alpha^k \text{ with } \sum_{j=0}^{m_k}\alpha^k_j = 1 \\
\quad x^{k+1} = \sum_{j=0}^{m_k}\alpha^k_j f(x^{k-m_k+j}) ,
\end{array}
\end{equation}
where $m$ is the memory (setting \texttt{acceleration\_lookback}, default
$10$; $0$ disables AA). Everything is in the choice of the weights. Define the
residual $g(x) = x - f(x)$, which vanishes at a fixed point, and let
$g_i = g(x^i)$ and
\[
Y_k = [\,y_{k-m_k}\ \cdots\ y_{k-1}\,],\quad y_i = g_{i+1} - g_i,
\qquad
S_k = [\,s_{k-m_k}\ \cdots\ s_{k-1}\,],\quad s_i = x^{i+1} - x^i .
\]

\paragraph{Type-II.} Choose the weights by least squares,
\[
\begin{array}{ll}
\minimize & \left\|\sum_{j=0}^{m_k} \alpha_j\, g(x^{k-m_k+j})\right\|_2^2\\[0.4em]
\mbox{subject to} & \sum_{j=0}^{m_k} \alpha_j = 1 .
\end{array}
\]
Reparametrizing by
$\gamma \in \R^{m_k}$ via $\alpha_0 = \gamma_0$,
$\alpha_i = \gamma_i - \gamma_{i-1}$, $\alpha_{m_k} = 1 - \gamma_{m_k-1}$ turns
this into $\minimize \|g_k - Y_k\gamma\|_2$, with solution
$\gamma^k = (Y_k\tr Y_k)^{-1}Y_k\tr g_k$ when $Y_k$ has full column rank. The
update becomes
\begin{equation}\label{eq:aa2}
x^{k+1} = x^k - B_k g_k, \qquad
B_k = I + (S_k - Y_k)(Y_k\tr Y_k)^{-1}Y_k\tr .
\end{equation}
$B_k$ minimizes $\|B - I\|_F$ subject to the inverse multi-secant condition
$BY_k = S_k$, so it is an approximate inverse Jacobian of $g$ and
\eqref{eq:aa2} is a quasi-Newton step --- $B_k$ is a generalized type-II
Broyden update of $I$.

\paragraph{Type-I.} Symmetrically, approximate the Jacobian itself, minimizing
$\|H - I\|_F$ subject to $HS_k = Y_k$, which gives
$H_k = I + (Y_k - S_k)(S_k\tr S_k)^{-1}S_k\tr$ and the update
$x^{k+1} = x^k - H_k^{-1}g_k$. The Woodbury identity yields the cheap form
\begin{equation}\label{eq:aa1}
H_k^{-1} = I + (S_k - Y_k)(S_k\tr Y_k)^{-1}S_k\tr ,
\end{equation}
in which the only inversion is of an $m_k \times m_k$ matrix.

Both types are available; type-I is the default
(\texttt{acceleration\_type\_1}) as it tends to perform better, while type-II
is more numerically stable but usually slower. Users switching to type-II
typically also lower \texttt{acceleration\_regularization}, whose default is
tuned for type-I.

\paragraph{Regularization.} A small $\epsilon > 0$ is added to the matrices
being inverted: $(Y_k\tr Y_k + \epsilon I)^{-1}$ for type-II, and
$(S_k\tr Y_k + \epsilon I)^{-1}$ for type-I (equivalent to regularizing
$S_k\tr S_k$ before applying Woodbury). Besides guaranteeing invertibility,
this shrinks the AA step towards the unaccelerated one: as
$\epsilon \to \infty$, $\gamma^\star \to 0$ and \eqref{eq:aa} reduces to
$x^{k+1} = f(x^k)$. Type-I typically needs more regularization than type-II.
The regularization folds into the factorizations by appending
$\sqrt{\epsilon}I$ below $S_k$ or $Y_k$.

\paragraph{In this implementation.} The linear algebra lives in
\texttt{src/aa.c}. The solve folds the Tikhonov rows into the
augmented matrix, takes its rank-revealing pivoted QR factorization with LAPACK
\texttt{geqp3}, truncates by rank, and applies a few steps of iterative
refinement to the $\gamma$ solve. This keeps the update stable as the columns of $S_k$ and $Y_k$
become nearly dependent near convergence. An SVD-based solve would be
similarly rank-revealing but is markedly more expensive per iteration.

The setting \texttt{acceleration\_interval} (default $5$) applies AA every
$k$th iteration, ignoring the intermediate ones. This makes AA $k$ times
cheaper, approximates a $k$ times longer memory, and decorrelates the data,
at the risk of staler iterates. Sweeps over the netlib LP and Maros--Meszaros
QP collections found $5$ robust: intervals near $1$ destabilize the
least-squares fit, and very large intervals leave acceleration idle. Solve
diagnostics --- acceptances, rejection causes, the rank of the last solve, the
last weight norm and the regularization used --- are reported in the
\texttt{aa\_stats} field of \texttt{ScsInfo}.

AA can dramatically speed up convergence, particularly when high accuracy is
wanted, and can also destabilize the solver badly. How best to deploy it
across all problems is an open question, and contributions in that direction
are welcome.

% ===========================================================================
\section{Termination}\label{sec:term}
% ===========================================================================

Because the Moreau decomposition \eqref{eq:moreau} is exact, the iterates
\emph{always} satisfy $s \in \K$, $y \in \K^*$ and $s \perp y$. Of the KKT
conditions \eqref{eq:kkt} only primal feasibility, dual feasibility and the
duality gap \eqref{eq:gap} can be violated, so termination reduces to three
residual tests. SCS stops and reports success when it has $x$, $s$, $y$ with
\begin{align}
r_p &:= \|Ax + s - b\|_\infty
      \leq \epsilon_{\rm abs} + \epsilon_{\rm rel}
            \max(\|Ax\|_\infty, \|s\|_\infty, \|b\|_\infty),
      \label{eq:tp}\\
r_d &:= \|Px + A\tr y + c\|_\infty
      \leq \epsilon_{\rm abs} + \epsilon_{\rm rel}
            \max(\|Px\|_\infty, \|A\tr y\|_\infty, \|c\|_\infty),
      \label{eq:td}\\
r_g &:= |x\tr Px + c\tr x + b\tr y|
      \leq \epsilon_{\rm abs} + \epsilon_{\rm rel}
            \max(|x\tr Px|, |c\tr x|, |b\tr y|),
      \label{eq:tg}
\end{align}
where $\epsilon_{\rm abs}$ and $\epsilon_{\rm rel}$ are the settings
\texttt{eps\_abs} and \texttt{eps\_rel} (both $10^{-4}$ by default). The
$\ell_\infty$ norm can be changed by redefining \texttt{NORM} in
\texttt{include/glbopts.h}.

Infeasibility is declared from the same iterates. SCS reports the problem
\textbf{primal infeasible} (dual unbounded) on finding $y$ with
\begin{equation}\label{eq:tpinf}
b\tr y = -1, \qquad
\|A\tr y\|_\infty < \epsilon_{\rm infeas}
  \cdot \min\!\left(1, \frac{\max_{ij}|A_{ij}|}{\|b\|_\infty}\right),
\end{equation}
and \textbf{dual infeasible} (primal unbounded) on finding $x$, $s$ with
\begin{equation}\label{eq:tdinf}
c\tr x = -1, \qquad
\|Px\|_\infty < \epsilon_{\rm infeas}, \qquad
\|Ax + s\|_\infty < \epsilon_{\rm infeas}
  \cdot \min\!\left(1, \frac{\max_{ij}|A_{ij}|}{\|c\|_\infty}\right),
\end{equation}
where $\epsilon_{\rm infeas}$ is the setting \texttt{eps\_infeas} (default
$10^{-7}$).

The $\min(1,\cdot)$ factors deserve a word. The normalizations $b\tr y = -1$
and $c\tr x = -1$ otherwise bake the magnitude of the data into the effective
tolerance: with $\|c\|_\infty \sim 10^6$, a normalized certificate $x$ has norm
$\sim 10^{-6}$, and the unscaled test is six orders of magnitude looser than
intended --- enough to declare a badly scaled feasible problem infeasible. The
factors only ever tighten the tests. As a further guard against transient
iterates that momentarily pass --- those produced by an AA step, for instance
--- the certificate conditions must hold at two consecutive residual checks
before SCS declares infeasibility.

% ===========================================================================
\section{The software}
% ===========================================================================

SCS is written in C99, and CI checks that it also compiles cleanly as
C++. It has no mandatory external dependencies: AMD and QDLDL are vendored
under \texttt{linsys/external}, and BLAS/LAPACK --- needed for the semidefinite
and spectral cones and for Anderson acceleration --- is on by default but can
be compiled out. The library is MIT licensed.

\subsection{Layout}

\begin{center}
\small
\begin{tabular}{@{}ll@{}}
\toprule
Path & Contents \\
\midrule
\texttt{include/}    & Public API (\texttt{scs.h}) and internal headers \\
\texttt{src/}        & Core solver: \texttt{scs.c}, cone projections, equilibration, AA, linear algebra \\
\texttt{src/spectral\_cones/} & Log-det, nuclear-norm, $\ell_1$ and sum-of-largest-eigenvalue cones \\
\texttt{linsys/}     & Linear-solver backends (Table~\ref{tab:linsys}) and matrix utilities \\
\texttt{linsys/external/} & Vendored AMD and QDLDL \\
\texttt{test/}       & Test suite (minunit), problem instances, packaging checks \\
\texttt{docs/src/}   & Sphinx documentation source \\
\texttt{docs/paper/} & This document \\
\bottomrule
\end{tabular}
\end{center}

\subsection{C API}

The one-shot entry point is
\begin{verbatim}
    scs_int scs(const ScsData *d, const ScsCone *k, const ScsSettings *stgs,
                ScsSolution *sol, ScsInfo *info);
\end{verbatim}
which returns an exit flag and fills \texttt{sol} with $(x, y, s)$ --- or with
a certificate --- and \texttt{info} with residuals, timings, iteration count
and AA diagnostics. Settings are initialized with
\texttt{scs\_set\_default\_settings}.

For a sequence of related problems the workspace is reused:
\begin{verbatim}
    ScsWork *w = scs_init(&d, &k, &stgs);     /* factorize, equilibrate */
    scs_solve(w, &sol, &info, /* warm_start = */ 0);
    scs_update(w, b_new, c_new);              /* new b and/or c */
    scs_solve(w, &sol, &info, /* warm_start = */ 1);
    scs_finish(w);
\end{verbatim}
As long as $A$ and $P$ are unchanged, \texttt{scs\_update} accepts new $b$ and
$c$ while the factorization and the equilibration stay cached, which is a large
saving. Warm starting is requested by the \texttt{warm\_start} setting, with
the guess supplied in the \texttt{x}, \texttt{y}, \texttt{s} members of
\texttt{ScsSolution} in the standard form \eqref{eq:problem}; those members are
overwritten with the solution at termination. Both mechanisms pay off when
solving a sequence of related problems --- a model-predictive-control loop, or
a full regularization path for the lasso.

\subsection{Build configuration}

Both a Makefile and a CMake build are provided:
\begin{verbatim}
    make                       # direct + indirect solvers, with BLAS/LAPACK
    make test && ./out/run_tests_direct

    cmake -DBUILD_TESTING=ON -DUSE_LAPACK=ON .. && cmake --build . && ctest
\end{verbatim}
The main compile-time options are:
\begin{center}
\small
\begin{tabular}{@{}lll@{}}
\toprule
Flag & Default & Effect \\
\midrule
\texttt{USE\_LAPACK}          & on  & BLAS/LAPACK; required for SD and spectral cones, and for AA \\
\texttt{DLONG}                & off & 64-bit integer indexing \\
\texttt{SFLOAT}               & off & Single-precision floats \\
\texttt{USE\_SPECTRAL\_CONES} & off & Spectral cones (experimental; requires LAPACK) \\
\texttt{USE\_OPENMP}          & off & OpenMP parallelism \\
\bottomrule
\end{tabular}
\end{center}
Each linear
solver is built into its own binary, so switching backends when linking
directly against SCS means linking the corresponding library; through a
language interface it is a setting.

\subsection{Interfaces}

The C library is wrapped for Python (\texttt{pip install scs}), Julia
(\texttt{SCS.jl}), R, MATLAB, Ruby, and JavaScript via WebAssembly. SCS is a
supported solver in CVX, CVXPY, YALMIP, Convex.jl and JuMP; in CVXPY it is the
default for semidefinite, exponential-cone and power-cone problems.

\subsection{Testing}

The suite in \texttt{test/} uses the minunit framework and runs against every
linear-solver backend that is built, checking solutions, infeasibility
certificates, warm starts, workspace reuse, settings round-tripping and
problem-file I/O. Spectral-cone tests are skipped unless the build enables
them. Continuous integration additionally covers the MKL, OpenMP, cuDSS and
GPU backends, CMake consumption by a downstream project, a Valgrind run, and
the documentation build.

% ===========================================================================
\section{Citing}\label{sec:citing}
% ===========================================================================

The original method and its initial numerical results are in
\cite{ocpb:16}; the quadratic extension and the homogeneous embedding
implemented here are in \cite{odonoghue:21}; the Anderson acceleration
scheme is from \cite{aa2020}. To cite a specific version of the software,
see \texttt{CITATION.cff} in the repository root.

% ===========================================================================
\begin{thebibliography}{9}

\bibitem{ocpb:16}
B.~O'Donoghue, E.~Chu, N.~Parikh and S.~Boyd,
\emph{Conic optimization via operator splitting and homogeneous self-dual
embedding},
Journal of Optimization Theory and Applications, 169(3):1042--1068, 2016.

\bibitem{odonoghue:21}
B.~O'Donoghue,
\emph{Operator splitting for a homogeneous embedding of the linear
complementarity problem},
SIAM Journal on Optimization, 31(3):1999--2023, 2021.

\bibitem{aa2020}
J.~Zhang, B.~O'Donoghue and S.~Boyd,
\emph{Globally convergent type-I Anderson acceleration for non-smooth
fixed-point iterations},
SIAM Journal on Optimization, 30(4):3170--3197, 2020.

\bibitem{anderson:65}
D.~G.~Anderson,
\emph{Iterative procedures for nonlinear integral equations},
Journal of the ACM, 12(4):547--560, 1965.

\bibitem{lions:79}
P.-L.~Lions and B.~Mercier,
\emph{Splitting algorithms for the sum of two nonlinear operators},
SIAM Journal on Numerical Analysis, 16(6):964--979, 1979.

\bibitem{boyd:04}
S.~Boyd and L.~Vandenberghe,
\emph{Convex Optimization},
Cambridge University Press, 2004.

\bibitem{ruiz:01}
D.~Ruiz,
\emph{A scaling algorithm to equilibrate both rows and columns norms in
matrices},
Technical Report RAL-TR-2001-034, Rutherford Appleton Laboratory, 2001.

\bibitem{amd}
P.~R.~Amestoy, T.~A.~Davis and I.~S.~Duff,
\emph{Algorithm 837: AMD, an approximate minimum degree ordering algorithm},
ACM Transactions on Mathematical Software, 30(3):381--388, 2004.

\bibitem{qdldl}
B.~Stellato, G.~Banjac, P.~Goulart, A.~Bemporad and S.~Boyd,
\emph{OSQP: an operator splitting solver for quadratic programs},
Mathematical Programming Computation, 12(4):637--672, 2020.

\bibitem{eigcg}
A.~Stathopoulos and K.~Orginos,
\emph{Computing and deflating eigenvalues while solving multiple right-hand
side linear systems with an application to quantum chromodynamics},
SIAM Journal on Scientific Computing, 32(1):439--462, 2010.

\end{thebibliography}

\end{document}
